% Generated by `board export` -- the transcript of a live lesson.
% Both halves of it: the tutor's cards and the student's own working.
%
% The three packages this file needs beyond the course's own are loaded AFTER
% coursemacros.sty and with no options, so that a course loading any of them
% its own way wins and there is no option clash -- which is a fatal error, and
% was one, the first time this was pointed at a real repository.
\PassOptionsToPackage{hidelinks}{hyperref}
\documentclass[11pt]{article}
% A scalable font, before anything else: a course whose macros switch microtype
% on gets font expansion, and font expansion against the default bitmap fonts is
% a fatal error rather than a warning. The course's own font package, if it has
% one, is loaded after this and wins.
\usepackage{lmodern}
\usepackage[T1]{fontenc}
\usepackage{amsmath,amssymb,amsthm,mathtools}
\usepackage[margin=1in]{geometry}
\usepackage{tikz}
% Commutative diagrams, if this machine has them. A minimal TeX install ships
% pgf but not tikz-cd, and requiring it unconditionally meant a missing package
% used by one optional kind of fence took down EVERY export on that machine:
% `! LaTeX Error: File 'tikz-cd.sty' not found`, on an evening's transcript that
% contained no diagrams at all. A lesson that cannot be exported is a lesson
% lost, so the whole document must not rest on a package it may not need.
\IfFileExists{tikz-cd.sty}{\usepackage{tikz-cd}}{}
\usetikzlibrary{arrows.meta,positioning,calc,fit,shapes.geometric}
\usepackage{coursemacros}
\input{board-macros.tex}
\usepackage{graphicx}
\usepackage{xcolor}
\usepackage{hyperref}
\usepackage{needspace}
% Who is speaking. A transcript in which the two voices look the same is a
% transcript nobody can read: the tutor's cards are the prose, and the
% student's own working is set off from it and labelled with whose it is.
\newenvironment{studentturn}[1]{%
  \par\medskip\needspace{4\baselineskip}%
  \noindent{\small\bfseries\sffamily #1}\par\nopagebreak\smallskip
  \begingroup\setlength{\leftskip}{1.2em}%
}{\endgroup\par\medskip}
\title{Ch 03 --- Rings}
\author{Mikey Ferguson}
\date{September 02, 2026}
\begin{document}
\maketitle
\needspace{3\baselineskip}\subsection*{Your move -- Chapter 3 --- the four we are doing, starting with 3.3}
Your list: 3.1, 3.2, \textbf{3.3} (circled), 3.8, 3.10, 3.11, 3.12 ("worth trouble"), 3.14, 3.15, 3.17, 3.25, 3.26. The bottom of the page is cut off, so whatever the arrow under 3.3 pointed at, I never saw it.

Thirteen is not a sitting. Four of them, in this order:

\begin{itemize}
  \item \textbf{3.3} --- circled, and it is the one that makes you \emph{use} "integral domain" instead of reciting it.
  \item \textbf{3.11} --- $a$ prime $\iff R/(a)$ an integral domain. The chapter's spine: ideals, quotients, primes in one line.
  \item \textbf{3.12} --- $\mathbb{Z}+i\sqrt{5}\,\mathbb{Z}$. Agreed, worth the trouble: it is \emph{the} object where irreducible $\neq$ prime.
  \item \textbf{3.15} --- roots and $(x-a)$, at most $n$ of them. Used in every chapter after this one.
\end{itemize}

Left out: 3.1, 3.2, 3.8, 3.10, 3.14, 3.17, 3.25, 3.26 --- deliberately, not forgotten. 3.14 and 3.25 come free once 3.11 lands.

\medskip\hrule\medskip

\textbf{Exercise 3.3.} \emph{Show that an integral domain with a finite number of elements is always a field.}

\begin{itemize}
  \item \emph{Ring:} commutative, with $1 \neq 0$.
  \item \emph{Integral domain:} a ring in which $rs = 0$ forces $r = 0$ or $s = 0$.
  \item \emph{Unit:} $q$ with $qq^{-1} = 1$ for some $q^{-1} \in R$.
  \item \emph{Field:} a ring in which every non-zero element is a unit.
\end{itemize}

Not yet. First, one thing by hand.

\textbf{Your move.} Take $a = 2$. Write out the full list of products $2x$ as $x$ runs over every element:

\begin{itemize}
  \item in $\mathbb{Z}_6$, for $x = 0,1,2,3,4,5$;
  \item in $\mathbb{Z}_5$, for $x = 0,1,2,3,4$.
\end{itemize}

Then answer two things about \emph{those lists}: in which one does $1$ appear, and in $\mathbb{Z}_6$ name a non-zero $x$ with $2x = 0$.


\begin{studentturn}{You said, 2026-09-02 10:35:38}
This is a homework set. We do them all.
\end{studentturn}

\needspace{3\baselineskip}\subsection*{Your move -- Homework --- all twelve, in the sheet's order. 3.1 first.}
Understood. Then nothing is chosen and nothing is dropped: twelve problems, worked in the sheet's order.

\textbf{3.1, 3.2, 3.3, 3.8, 3.10, 3.11, 3.12, 3.14, 3.15, 3.17, 3.25, 3.26.}

Skip any one and I leave it and move on, then bring it back at the end --- a skipped problem is an empty solution region, not a deleted one. The $\mathbb{Z}_6$ / $\mathbb{Z}_5$ check from the last card is parked; it comes back when we reach 3.3.

\medskip\hrule\medskip

\textbf{Problem 3.1.} *Suppose that $S$ is a set and $R$ is a ring. Let $R^S$ denote the set of all mappings from $S$ to $R$. Show that $R^S$ is a ring, under the operations defined by*
$$(f+g)(s) = f(s) + g(s), \qquad (fg)(s) = f(s)g(s).$$
*Show that if $S$ has more than one element then there exist non-zero elements $f$ and $g$ in $R^S$ for which $fg = 0$.*

\begin{itemize}
  \item \emph{Ring:} commutative under both operations, associative, distributive, with an identity $1 \neq 0$; abelian group under addition.
  \item *$R^S$:* its elements are \textbf{functions} $S \to R$. Two are equal iff they agree at every $s$.
  \item *Non-zero in $R^S$:* not the zero element of $R^S$ --- so it takes a non-zero value somewhere, not everywhere.
  \item \emph{Integral domain:} $rs = 0 \Rightarrow r = 0$ or $s = 0$. The second half exhibits $R^S$ failing this.
\end{itemize}

Not yet. One thing by hand first.

\textbf{Your move.} Let $R = \mathbb{Z}$ and $S = \{a, b\}$, and write a function as its pair of values. Take $f(a) = 1$, $f(b) = 0$ and $g(a) = 0$, $g(b) = 1$.

Compute $(fg)(a)$ and $(fg)(b)$, and then say what the zero element of $R^S$ actually is and whether $f$ equals it.


\begin{studentturn}{You wrote -- attempt 1 of 4, 2026-09-02 10:59:05}
\begin{center}\includegraphics[width=0.95\linewidth,height=0.78\textheight,keepaspectratio]{live/answers/t0002-r1.png}\end{center}
\end{studentturn}

\begin{studentturn}{You wrote -- attempt 2 of 4, 2026-09-02 11:05:20}
\begin{center}\includegraphics[width=0.95\linewidth,height=0.78\textheight,keepaspectratio]{live/answers/t0002-r2.png}\end{center}
\end{studentturn}

\begin{studentturn}{You wrote -- attempt 3 of 4, 2026-09-02 11:13:53}
\begin{center}\includegraphics[width=0.95\linewidth,height=0.78\textheight,keepaspectratio]{live/answers/t0002-r3.png}\end{center}
\end{studentturn}

\begin{studentturn}{You wrote -- attempt 4 of 4, 2026-09-02 11:16:03}
\begin{center}\includegraphics[width=0.95\linewidth,height=0.78\textheight,keepaspectratio]{live/answers/t0002-r4.png}\end{center}
\end{studentturn}

\needspace{3\baselineskip}\subsection*{Not quite -- The pair is right. "Clearly" is not a proof.}
Settled, and I will not ask for it again:

\begin{itemize}
  \item The zero of $R^S$ is the constant function $s \mapsto 0_R$ and the identity is $s \mapsto 1_R$. Both correct.
  \item The pair in the lower block is exactly the right pair: $f$ equal to $1_R$ at $s_1$ and $0_R$ elsewhere, $g$ equal to $1_R$ at $s_2$ and $0_R$ elsewhere, with $s_1 \neq s_2$. Their product is the zero function, and for the right reason --- at every $s$, at least one of the two factors is $0_R$.
\end{itemize}

\textbf{The break is the third line of the upper block.} \emph{"Clearly by the specified function properties we have all ring properties"} is the thing 3.1 asks you to show, restated as though it had been shown. It names no axiom, and it never uses the one fact the whole problem turns on: both operations are defined \textbf{pointwise}, so an identity in $R^S$ is the corresponding identity in $R$, evaluated at each $s$ separately.

\textbf{Two smaller things in the lower block.} What you exhibited is called a \textbf{zero divisor} --- a non-zero element whose product with another non-zero element is $0$. "Non-zero divisor" is the opposite property. And nothing there says \emph{why} $f$ and $g$ are non-zero; that is one line, and it needs an axiom of $R$.

\textbf{Your move.} One axiom, written out in full, to fix the shape. Let $f, g, h \in R^S$ and prove left distributivity:
$$f(g+h) = fg + fh.$$
Two elements of $R^S$ are equal iff they agree at every $s \in S$. So fix an arbitrary $s \in S$, evaluate each side at $s$ using only the two defining formulas, and name the property of $R$ that closes the gap between them.

Four or five lines. This one axiom, not all of them.


\needspace{3\baselineskip}\subsection*{Not quite -- f(h(s)) is not a thing that exists}
"Zero divisors" is now right. The lower block still never says why $f$ and $g$ are non-zero --- one line, still owed.

\textbf{The break is the second line of the new block.} You wrote

$$(f+g)(h(s)) = f(h(s)) + g(h(s)).$$

Look at what $f$ eats. Elements of $R^S$ are functions $S \to R$: $f$ takes an element of $S$ and returns an element of $R$. But $h(s)$ is an element of $R$. So $f(h(s))$ feeds a ring element into a function that is only defined on $S$ --- it does not typecheck, and no rearrangement of that line will fix it.

What you have written is \textbf{composition} of functions. The multiplication in $R^S$ is not composition. Go back to the defining formula you were given: $(fg)(s) = f(s)g(s)$ --- the right-hand side is a *product of two elements of $R$*, not one function fed into another. Nothing in this problem ever puts a function inside a function.

\textbf{And the identity you wrote down is false as stated.} $(f+g)h$ does not pair with $fg + fh$. Multiply out $(a+b)c$ in any ring and read off which two products you get.

\textbf{Your move.} Same job, stated correctly this time. Let $f, g, h \in R^S$ and prove

$$(f+g)h = fh + gh.$$

Fix an arbitrary $s \in S$ and work at that $s$. Every line you write must be an equation between **elements of $R$**. You are allowed exactly two tools: the two defining formulas, and one named axiom of $R$ --- and I want that axiom named on the page where you use it.

Four or five lines, every one of them an equality in $R$.


\needspace{3\baselineskip}\subsection*{Not quite -- The first line finally lands. The last one is false.}
\textbf{This landed, and it is the step you were missing for two revisions:}

$$\big((f+g)h\big)(s) = \big(f(s) + g(s)\big)\,h(s).$$

You evaluated at $s$ and came out with a product of two elements of $R$. No function sits inside another function anywhere in it. That is the move the whole problem runs on, and it is now on the page.

\textbf{The break is the last equality.} You wrote

$$f(s)h(s) + g(s)h(s) \;=\; fg(s) \cdot gh(s).$$

A sum went in and a product came out. Nothing between those two expressions can turn a $+$ into a $\cdot$, and the identity you are proving has a $+$ on its right-hand side --- so the chain has to end with one.

The second factor is wrong too, and the same slip is in the statement you copied at the top: you have written $fg$ where the identity needs $fh$. Multiply out $(a+b)c$ once more and read off which two products it is.

\textbf{Two notes on the middle line.} $f(s)$ and $h(s)$ are both elements of $R$, so they are multiplied by $R$'s own multiplication --- $\circ$ means composition, and there is none in this problem. And I asked for the axiom of $R$ you used to be named at the point you use it. It is still not there.

\textbf{Your move.} Start from

$$f(s)h(s) + g(s)h(s)$$

and finish the chain, ending at $(fh + gh)(s)$. Every step an equality in $R$.

The one thing this needs: a defining formula is an \emph{equation}, so it reads in both directions. You have been using $(fh)(s) = f(s)h(s)$ left to right all along. The end of this chain needs it right to left.

Still owed, third time: one line saying why $f$ and $g$ are non-zero.


\needspace{3\baselineskip}\subsection*{Your move -- Down a level --- three integers}
The $\cdot$ is now a $+$, and the tail of the chain collapses correctly into a single function evaluated at $s$. Both settled.

One thing has not moved in three revisions, and I have flagged it twice in passing, which was evidently the wrong way to flag it. Your statement, top right, reads

$$(f+g)h = fg + gh,$$

and the last line of your chain ends the same way. That pattern --- first two letters, then last two --- is not the distributive law. Nothing about $R^S$ is involved in this; it is going wrong in $R$.

So we leave functions entirely for one minute and do it where you can simply compute.

\textbf{Your move.} Take $a = 2$, $b = 3$, $c = 4$ in $\mathbb{Z}$. Work out all three of these as plain integers:

$$(a+b)c, \qquad ab + bc, \qquad ac + bc.$$

Then say which of the two right-hand expressions equals the left-hand one --- and which of the two has the shape you have been writing on your page.

Parked, not asked about this turn, and coming back straight after it: the $\circ$ sitting between two elements of $R$, the axiom of $R$ you have still not named, and the one line saying why $f$ and $g$ are non-zero.


\begin{studentturn}{You wrote, 2026-09-02 11:18:29}
\begin{center}\includegraphics[width=0.95\linewidth,height=0.78\textheight,keepaspectratio]{live/answers/t0003-r1.png}\end{center}
\end{studentturn}

\needspace{3\baselineskip}\subsection*{Your move -- The chain is right. The heading above it is not.}
$fh(s) + gh(s)$. That is the correct end of the chain, and the line is now sound from $\big((f+g)h\big)(s)$ all the way down. Settled.

\textbf{Bookkeeping, no thinking required.} The two lines at the top of that block still say $fg + gh$ --- *"Prove $(f+g)h = fg + gh$"\emph{ and }"show $\big((f+g)h\big)(s) = fg(s) + gh(s)$"*. Your chain no longer ends there. As the page stands it announces one identity and proves a different one. Change both to $fh$.

\textbf{Two notes on notation.} Write $(fh)(s)$ rather than $fh(s)$: the bracket says *the function $fh$, evaluated at $s$*, and that distinction is the entire problem. And the $\circ$ between $f(s)$ and $h(s)$ should be $R$'s own multiplication --- there is no composition anywhere in 3.1.

\textbf{Your move.} One line, and it is the only piece of mathematical content in the whole first half of 3.1. Look at this step of your chain:

$$\big(f(s) + g(s)\big)\,h(s) \;=\; f(s)h(s) + g(s)h(s).$$

Every quantity in it is an element of $R$, and the equality is not something you proved --- it is something you were given. Name the axiom of $R$ that supplies it, and write that name on the page beside the step.

That is what "$R^S$ is a ring" actually rests on, and it is why one axiom was enough to ask for: the other axioms go the same way, each one handed over by its counterpart in $R$.

Then the non-zero line for $f$ and $g$, on its own. Then 3.1 in full.


\begin{studentturn}{You wrote, 2026-09-02 11:20:23}
\begin{center}\includegraphics[width=0.95\linewidth,height=0.78\textheight,keepaspectratio]{live/answers/t0004-r1.png}\end{center}
\end{studentturn}

\needspace{3\baselineskip}\subsection*{Your move -- Problem 3.1, in full}
Headings fixed. The distributivity line is complete and correct, and that was the whole point of asking for one axiom: each axiom of $R^S$ is its counterpart in $R$, evaluated at a fixed arbitrary $s$. The rest go the same way, so 3.1 goes back to you whole.

\textbf{Problem 3.1.} *Suppose that $S$ is a set and $R$ is a ring. Let $R^S$ denote the set of all mappings from $S$ to $R$. Show that $R^S$ is a ring, under the operations defined by*

$$(f+g)(s) = f(s) + g(s), \qquad (fg)(s) = f(s)g(s).$$

*Show that if $S$ has more than one element then there exist non-zero elements $f$ and $g$ in $R^S$ for which $fg = 0$.*

Everything the statement leans on:

\begin{itemize}
  \item \emph{Ring:} abelian group under $+$; multiplication associative and commutative with an identity $1 \neq 0$; multiplication distributes over $+$.
  \item *$R^S$:* its elements are \textbf{functions} $S \to R$. Two are equal iff they agree at every $s \in S$.
  \item *Zero and identity of $R^S$:* the constant functions $s \mapsto 0_R$ and $s \mapsto 1_R$. You have both.
  \item *Non-zero in $R^S$:* **not equal to the zero element of $R^S$** --- so it takes a non-zero value at \emph{some} $s$, not at every $s$.
  \item *$fg = 0$ in $R^S$:* the function $fg$ \emph{is} the zero function, i.e. $(fg)(s) = 0_R$ for \textbf{every} $s \in S$.
\end{itemize}

\textbf{Your move.} Write 3.1 out as a finished answer. Keep what you have --- the distributive axiom and the pair $f, g$ --- and do not redo them.

Three things are not on your page yet: the remaining ring axioms, by the same pointwise argument; the reason each of $f$ and $g$ is non-zero; and the check that $(fg)(s) = 0_R$ at \emph{every} $s$, not only at $s_1$ and $s_2$.


\begin{studentturn}{You said, 2026-09-02 11:22:26}
What I wrote is good enough. Type it up and we're moving on.
\end{studentturn}

\needspace{3\baselineskip}\subsection*{Your move -- 3.1 typed up as written. Now 3.2.}
Going into 3.1's solution region exactly as you wrote it: the identification of $0_{R^S}$ and $1_{R^S}$, the distributivity chain, and the zero-divisor pair. Your steps, your order, nothing added.

So that you know what is in the file rather than finding out later: it establishes \textbf{one} of the ring axioms for $R^S$, and the second half exhibits $f$ and $g$ without the line asserting they are non-zero. The region stays yours --- add to it whenever you like and I will re-typeset in seconds.

Moving on.

\textbf{Problem 3.2.} *Suppose that $R$ is an integral domain, with field of fractions $F$. Show that the field of fractions of $R[x_1, \ldots, x_n]$ can be identified naturally with $F(x_1, \ldots, x_n)$.*

Everything the statement uses:

\begin{itemize}
  \item \emph{Integral domain:} a commutative ring with $1 \neq 0$ in which $rs = 0$ implies $r = 0$ or $s = 0$.
  \item *Field of fractions of an integral domain $D$:* the field of formal quotients $a/b$ with $a, b \in D$ and $b \neq 0$, where $a/b = c/d$ exactly when $ad = bc$. $D$ sits inside it by $a \mapsto a/1$.
  \item *$R[x_1,\ldots,x_n]$:* polynomials in $n$ indeterminates with coefficients in $R$ --- itself an integral domain.
  \item *$F(x_1,\ldots,x_n)$:* by definition the field of fractions of $F[x_1,\ldots,x_n]$.
  \item \emph{Identified naturally:} exhibit an isomorphism between the two fields that fixes $R$ and each $x_i$.
\end{itemize}

\textbf{Your move.} 3.2, cold. Ask if you want it broken down; otherwise I will read what you write.


\begin{studentturn}{You wrote -- attempt 1 of 17, 2026-09-02 11:34:39}
\begin{center}\includegraphics[width=0.95\linewidth,height=0.78\textheight,keepaspectratio]{live/answers/t0006-r1.png}\end{center}
\end{studentturn}

\begin{studentturn}{You wrote -- attempt 2 of 17, 2026-09-02 11:52:19}
\begin{center}\includegraphics[width=0.95\linewidth,height=0.78\textheight,keepaspectratio]{live/answers/t0006-r2.png}\end{center}
\end{studentturn}

\begin{studentturn}{You wrote -- attempt 3 of 17, 2026-09-02 12:19:43}
\begin{center}\includegraphics[width=0.95\linewidth,height=0.78\textheight,keepaspectratio]{live/answers/t0006-r3.png}\end{center}
\end{studentturn}

\begin{studentturn}{You wrote -- attempt 4 of 17, 2026-09-02 12:42:34}
\begin{center}\includegraphics[width=0.95\linewidth,height=0.78\textheight,keepaspectratio]{live/answers/t0006-r4.png}\end{center}
\end{studentturn}

\begin{studentturn}{You wrote -- attempt 5 of 17, 2026-09-02 13:01:48}
\begin{center}\includegraphics[width=0.95\linewidth,height=0.78\textheight,keepaspectratio]{live/answers/t0006-r5.png}\end{center}
\end{studentturn}

\begin{studentturn}{You wrote -- attempt 6 of 17, 2026-09-02 13:04:57}
\begin{center}\includegraphics[width=0.95\linewidth,height=0.78\textheight,keepaspectratio]{live/answers/t0006-r6.png}\end{center}
\end{studentturn}

\begin{studentturn}{You wrote -- attempt 7 of 17, 2026-09-02 13:06:56}
\begin{center}\includegraphics[width=0.95\linewidth,height=0.78\textheight,keepaspectratio]{live/answers/t0006-r7.png}\end{center}
\end{studentturn}

\begin{studentturn}{You wrote -- attempt 8 of 17, 2026-09-02 13:09:07}
\begin{center}\includegraphics[width=0.95\linewidth,height=0.78\textheight,keepaspectratio]{live/answers/t0006-r8.png}\end{center}
\end{studentturn}

\begin{studentturn}{You said -- attempt 9 of 17, 2026-09-02 13:22:13}
Remind me of the definition of isomorphism again, and maybe then I'll understand why you're saying I need to be careful about what I'm mapping to what, because I was a under the impression I was mapping R[x1\ldots{}xn] to F(x1\ldots{}xn)\ldots{}
\end{studentturn}

\begin{studentturn}{You wrote -- attempt 10 of 17, 2026-09-02 13:31:46}
\begin{center}\includegraphics[width=0.95\linewidth,height=0.78\textheight,keepaspectratio]{live/answers/t0006-r9.png}\end{center}
\end{studentturn}

\begin{studentturn}{You wrote -- attempt 11 of 17, 2026-09-02 13:36:01}
\begin{center}\includegraphics[width=0.95\linewidth,height=0.78\textheight,keepaspectratio]{live/answers/t0006-r10.png}\end{center}
\end{studentturn}

\begin{studentturn}{You said -- attempt 12 of 17, 2026-09-02 13:41:10}
I don't understand
\end{studentturn}

\begin{studentturn}{You wrote -- attempt 13 of 17, 2026-09-02 13:44:46}
\begin{center}\includegraphics[width=0.95\linewidth,height=0.78\textheight,keepaspectratio]{live/answers/t0006-r11.png}\end{center}
\end{studentturn}

\begin{studentturn}{You wrote -- attempt 14 of 17, 2026-09-02 15:33:57}
\begin{center}\includegraphics[width=0.95\linewidth,height=0.78\textheight,keepaspectratio]{live/answers/t0006-r12.png}\end{center}
\end{studentturn}

\begin{studentturn}{You wrote -- attempt 15 of 17, 2026-09-02 15:39:28}
\begin{center}\includegraphics[width=0.95\linewidth,height=0.78\textheight,keepaspectratio]{live/answers/t0006-r13.png}\end{center}
\end{studentturn}

\begin{studentturn}{You wrote -- attempt 16 of 17, 2026-09-02 15:40:46}
\begin{center}\includegraphics[width=0.95\linewidth,height=0.78\textheight,keepaspectratio]{live/answers/t0006-r14.png}\end{center}
\end{studentturn}

\begin{studentturn}{You said -- attempt 17 of 17, 2026-09-02 15:43:24}
2
\end{studentturn}

\needspace{3\baselineskip}\subsection*{Aside -- What $x_1,\ldots,x_n$ is, and one of each creature}
Fair question. The notation is doing nothing clever.

**$x_1, \ldots, x_n$ is just $n$ different variables.** The subscript is a label --- not a power, not a value. $n$ is any fixed positive integer, and nothing in 3.2 depends on which one, so take $n = 2$ for the whole problem and write $x_1, x_2$. Read them as $x$ and $y$ if that is easier.

For concrete creatures, take $R = \mathbb{Z}$. Then $F = \mathbb{Q}$.

**The left-hand one: $\mathbb{Z}[x_1, x_2]$.** Polynomials in two variables with integer coefficients. For instance

$$3x_1^2x_2 - 5x_2 + 7.$$

Your brace is right --- every element is a polynomial. Note what that \emph{excludes}: $\tfrac{1}{2}x_1$ is not in it, because $\tfrac{1}{2} \notin \mathbb{Z}$.

**The right-hand one: $\mathbb{Q}(x_1, x_2)$.** Quotients of polynomials with rational coefficients. For instance

$$\frac{x_1^2 + \tfrac{1}{2}x_2}{x_1x_2 - 3}.$$

Your other brace is right too, and here is the equivalence actually biting:

$$\frac{x_1x_2}{x_1} = \frac{x_2}{1}, \qquad \text{because } x_1x_2 \cdot 1 = x_1 \cdot x_2.$$

One element, two names. That is what "equivalence class $a/b$" buys you.

\textbf{Your move.} Three expressions. For each one, say whether it lies in $\mathbb{Z}[x_1,x_2]$, in $\mathbb{Q}(x_1,x_2)$, or in both:

\begin{enumerate}
  \item $4x_1^3 - x_1x_2 + 2$
  \item $\tfrac{1}{3}x_2 + 5$
  \item $\dfrac{1}{x_1 + x_2}$
\end{enumerate}


\needspace{3\baselineskip}\subsection*{Aside -- Mostly right. The arrow starts in the wrong place.}
Asked and answered, item by item.

\textbf{The sorting: all three correct.} $4x_1^3 - x_1x_2 + 2$ in both, $\tfrac13x_2 + 5$ only in $\mathbb{Q}(x_1,x_2)$, $1/(x_1+x_2)$ only in $\mathbb{Q}(x_1,x_2)$. Nothing to fix.

**$a, b \in F[x_1,\ldots,x_n]$: right.** Add the one condition you left off --- $b \neq 0$, or there is no fraction.

**"$R^n \to R$ with coefficients in $R$": not right.** A polynomial is a \emph{formal expression} --- a finite list of coefficients attached to monomials. It \emph{gives} you a function on $R^n$, but it is not that function, and the two are genuinely different objects: over $\mathbb{F}_2$, $x^2$ and $x$ take the same value at every point and are still different polynomials. You proved 3.1 about a ring of functions; $R[x_1,\ldots,x_n]$ is not one.

\textbf{And now the thing your bubble is really asking about.} Look at where your $\phi$ starts:

$$\phi : R[x_1,\ldots,x_n] \longrightarrow F(x_1,\ldots,x_n).$$

3.2 does not ask you to identify $R[x_1,\ldots,x_n]$ with $F(x_1,\ldots,x_n)$. It asks you to identify \textbf{the field of fractions of} $R[x_1,\ldots,x_n]$ with $F(x_1,\ldots,x_n)$. Those are not the same object, and no map from the left-hand set as you have written it can be an isomorphism onto the right.

\textbf{Your move.} Take $R = \mathbb{Z}$ and $n = 2$, so the right-hand side is $\mathbb{Q}(x_1,x_2)$.

Does $x_1$ have a multiplicative inverse in $\mathbb{Z}[x_1,x_2]$? Does it have one in $\mathbb{Q}(x_1,x_2)$? Then write down what the domain of $\phi$ has to be.


\needspace{3\baselineskip}\subsection*{Aside -- No. A polynomial is not the function it gives you.}
Direct answer to the bubble: \textbf{no.}

A polynomial in $R[x_1,\ldots,x_n]$ is a \emph{formal} object --- a finite list of coefficients from $R$, one per monomial. Two polynomials are equal exactly when their coefficients agree term by term. Nothing about values enters into it.

Every polynomial does \emph{hand you} a function $R^n \to R$, by substituting elements of $R$ for the variables. But that is something a polynomial produces, not what it is, and the difference is not pedantry: the passage from polynomial to function loses information.

Here is the loss, concretely. Take $R = \mathbb{F}_2 = \{0,1\}$ and $n = 1$, and consider

$$p = x^2, \qquad q = x.$$

As functions on $\mathbb{F}_2$: $p(0) = 0 = q(0)$ and $p(1) = 1 = q(1)$ --- \textbf{the same function}. As polynomials they are plainly different: different coefficients, different degrees. So the functions $R^n \to R$ form a coarser collection than $R[x_1,\ldots,x_n]$, and swapping one for the other changes the object.

It matters for the thing you got right below, too. $x_1$ has no inverse in $\mathbb{Z}[x_1,x_2]$ because \textbf{degrees add} --- the degree of a product is the sum of the degrees, so nothing of degree $1$ can multiply up to the constant $1$. That is a fact about coefficients. It is invisible in the values.

\textbf{Both inverse answers are correct}, and $1/x_1$ is the right witness. You do not need the caveat: $x_1$ is an indeterminate, so it is never $0_R$. Still not written down: **what the domain of $\phi$ has to be.** That one is owed.

One flat correction. Your right-hand line now reads $a, b \in R[x_1,\ldots,x_n]$; last revision it read $F[x_1,\ldots,x_n]$, which was right. $F(x_1,\ldots,x_n)$ is \emph{by definition} the field of fractions of $F[x_1,\ldots,x_n]$. Put the $F$ back.

\textbf{Your move.} Stay in $\mathbb{F}_2[x]$, one variable. Write down a polynomial that is \textbf{not} the zero polynomial but whose function sends every element of $\mathbb{F}_2$ to $0$.


\needspace{3\baselineskip}\subsection*{Aside -- Yes. $F$ is not some other field --- it is built out of $R$.}
Direct answer to the bubble: \textbf{yes}, and it is the tightest relation available.

The hypothesis of 3.2 says it outright: *$R$ is an integral domain, with field of fractions $F$.* $F$ is not a field handed to you alongside $R$. It is manufactured from $R$ and nothing else:

$$F = \{\, a/b \;:\; a,b \in R,\; b \neq 0_R \,\}, \qquad a/b = c/d \iff ad = bc.$$

Every symbol on the right-hand side of that is an element of $R$.

And $R$ sits \emph{inside} $F$: the map $a \mapsto a/1_R$ is an injective ring homomorphism, so $F$ contains a faithful copy of $R$. That is the whole content of "relate".

Your example is right for exactly this reason: $\mathbb{Q}$ \emph{is} the field of fractions of $\mathbb{Z}$. Had you written $R = \mathbb{Z}$, $F = \mathbb{R}$, the hypothesis would simply be false --- $\mathbb{R}$ is a field containing $\mathbb{Z}$, but it is not $\mathbb{Z}$'s field of fractions. Containing $R$ is not enough; $F$ has to be \emph{generated} by quotients of elements of $R$.

Why this is the hinge of 3.2 and not a side remark: both objects in the statement are towers built on the same ground floor. On the left, the field of fractions of $R[x_1,\ldots,x_n]$. On the right, $F(x_1,\ldots,x_n)$ --- which is the field of fractions of $F[x_1,\ldots,x_n]$, and $F$ itself came from $R$. One ring at the bottom of both. That shared bottom is what makes an identification \emph{natural} rather than a coincidence.

\textbf{Your move.} Take $R = \mathbb{Z}$, $F = \mathbb{Q}$, with the embedding $a \mapsto a/1_R$. Write down which element of $\mathbb{Q}$ the integer $3$ is sent to, and then name one element of $\mathbb{Q}$ that is \textbf{not} in the image of that map.


\needspace{3\baselineskip}\subsection*{Not quite -- $3 \mapsto 3/1$ is right. "$0$ is not mapped" is false.}
The equivalence relation line is correct as written, and $3 \mapsto 3/1$ is correct.

The last line is where it breaks. **$0_R$ is mapped.** The embedding $a \mapsto a/1_R$ is defined for \emph{every} $a \in R$ with no exception, so $0 \mapsto 0/1_R$, and $0/1_R$ is $0_F$. Nothing in $R$ is left out --- that is what makes it a ring homomorphism at all: it has to send $0_R$ to $0_F$ and $1_R$ to $1_F$.

Where $0_R$ is genuinely forbidden is the \textbf{denominator}, and that is a different slot in the same expression. Which is the flat correction on the line above it: you wrote

$$F = \{\, r_1/r_2 \;:\; r_1, r_2 \in R \,\}$$

and it needs $r_2 \neq 0_R$. Numerators range over all of $R$, including $0_R$; denominators over $R \setminus \{0_R\}$. You already had that condition written correctly on the right-hand side of the diagram for $b$.

So the question is still open, and it is asking for something else. The image of $\mathbb{Z} \to \mathbb{Q}$ is the set of all fractions of the form $a/1$ --- a faithful copy of $\mathbb{Z}$ sitting inside $\mathbb{Q}$. Every integer is in there. But $\mathbb{Q}$ is \emph{strictly bigger} than that copy, and the elements it has left over are the entire reason $F$ is worth constructing.

\textbf{Your move.} Name one element of $\mathbb{Q}$ that is \textbf{not} equal to $a/1$ for any integer $a$ --- and say in one line why no integer maps to it.


\needspace{3\baselineskip}\subsection*{Not quite -- $4/2$ has $b \neq 1_R$. And $2$ maps to it.}
You dropped "$0$ is not mapped", which was the right move. The replacement is still false, and it fails for a reason worth having.

You wrote: *nothing maps to any $a/b$ with $b \neq 1_R$.*

Take $4/2$. Its denominator is $2 \neq 1_R$, so your claim says no integer maps to it. But $2 \mapsto 2/1$, and $2/1 = 4/2$ --- check it against the relation you wrote one line up: $2 \cdot 2 = 4 \cdot 1$. So $4/2$ \textbf{is} in the image, and it is one of your own excluded ones.

The break is not arithmetic. It is that "$b \neq 1_R$" is a property of a \textbf{representative}, not of the element. Elements of $F$ are equivalence classes; a single class has infinitely many names, $2/1$ and $4/2$ and $-6/{-3}$, and they differ in what $b$ is. So no condition written on $b$ alone can pick out a set of elements --- you have described a property of \emph{how something is written}, and the thing you want to describe is what it \emph{is}.

Which is why the question asked for \textbf{one element}, not a rule. A single concrete rational, where you can say plainly that \emph{no} integer maps to it --- that is a claim about every way of writing it, and it will force you to say the right thing.

\textbf{Your move.} Name one specific element of $\mathbb{Q}$ that no integer maps to, and give the one-line reason.

Still outstanding from the last card, untouched: your line $F = \{r_1/r_2 : r_1, r_2 \in R\}$ needs $r_2 \neq 0_R$.


\needspace{3\baselineskip}\subsection*{Not quite -- $\gcd(a,b)=1$ fixes the real problem. $a \neq b$ is the wrong condition.}
The $\gcd(a,b) = 1$ is a genuine repair and you should keep it. Every non-zero rational has exactly one representative in lowest terms, so "the $a/b$ with $\gcd(a,b)=1$" names a \emph{canonical} representative --- which makes a condition written on it a property of the element after all. That was the objection from last card and it is now answered.

The second clause is where it breaks. Take $3/1$. $\gcd(3,1) = 1$, and $a = 3$, $b = 1$, so $a \neq b$. Your rule therefore claims nothing maps to $3/1$ --- two lines above it, on your own page, $3 \mapsto 3/1$.

Look at what $a \neq b$ actually removes. In lowest terms, $a = b$ forces $\gcd(a,a) = a = 1$, so the only class it excludes is $1/1$. One rational, and it is an integer, so it is not even the kind of thing you were trying to exclude. The clause is doing almost nothing and what little it does is backwards.

The set you are reaching for is the integers-inside-$\mathbb{Q}$, and in lowest terms there is a clean way to say which classes those are. It is not a comparison between $a$ and $b$.

\textbf{Your move.} Rewrite the second clause so that your rule rejects $3/1$ --- and $-7/1$, and every other integer --- while still keeping the rationals that are not integers. One line.


\needspace{3\baselineskip}\subsection*{Correct -- $b \neq 1$ is right. Now say the same thing one level up.}
Correct. In lowest terms an element of $\mathbb{Q}$ is an integer exactly when its denominator is $1$, so the classes nothing maps to are exactly those with $\gcd(a,b)=1$ and $b \neq 1$. The step that carried it was keeping $\gcd(a,b)=1$: it pins a single representative, and only then does a condition on $b$ mean anything.

So $\mathbb{Z} \hookrightarrow \mathbb{Q}$ is injective and \textbf{not} surjective, and the things it misses are the whole point of $\mathbb{Q}$.

Now look at your arrow, which has not changed in eight revisions:

$$\phi : R[x_1,\ldots,x_n] \longrightarrow F(x_1,\ldots,x_n).$$

That is the same shape of map as $\mathbb{Z} \to \mathbb{Q}$, one level up --- and you already settled, earlier, that $x_1$ has no inverse in $\mathbb{Z}[x_1,x_2]$ while $1/x_1$ sits in $\mathbb{Q}(x_1,x_2)$. So $\phi$ misses things, exactly as $\mathbb{Z} \to \mathbb{Q}$ does.

\textbf{3.2, in full:} *Suppose that $R$ is an integral domain, with field of fractions $F$. Show that the field of fractions of $R[x_1,\ldots,x_n]$ can be identified naturally with $F(x_1,\ldots,x_n)$.*

\begin{itemize}
  \item *Field of fractions of $D$:* the classes $a/b$, $a,b \in D$, $b \neq 0_D$, with $a/b = c/d \iff ad = bc$.
  \item *$F(x_1,\ldots,x_n)$:* by definition the field of fractions of $F[x_1,\ldots,x_n]$.
  \item \emph{Identified:} an isomorphism of rings, not merely an injection.
\end{itemize}

\textbf{Your move.} Write the corrected arrow: what does the domain of $\phi$ have to be?


\needspace{3\baselineskip}\subsection*{Aside -- Isomorphism, and why a non-field cannot be one with a field}
\textbf{Definition.} A map $\theta : A \to B$ between rings is a \emph{ring isomorphism} if it is

\begin{enumerate}
  \item a ring homomorphism: $\theta(a + a') = \theta(a) + \theta(a')$, $\theta(aa') = \theta(a)\theta(a')$, and $\theta(1_A) = 1_B$; and
  \item a \textbf{bijection} --- injective \emph{and} surjective.
\end{enumerate}

Clause 2 is the whole of your problem. "Can be identified naturally with" in 3.2 means \emph{there is an isomorphism}, and "naturally" only adds that it is the canonical one rather than some arbitrary choice.

Your impression is half right, and the right half is worth keeping: there \textbf{is} a map $R[x_1,\ldots,x_n] \to F(x_1,\ldots,x_n)$, it is a homomorphism, and it is injective. It is the natural place to start building. It is simply not the isomorphism, and it cannot be made into one.

Here is why, and you already have every piece. An isomorphism transports structure in both directions: if $A \cong B$ and $B$ is a field, then $A$ is a field, because you can carry an inverse back along $\theta^{-1}$. Now --- $F(x_1,\ldots,x_n)$ is a field. $R[x_1,\ldots,x_n]$ is not: $x_1$ has no inverse in $\mathbb{Z}[x_1,x_2]$, as you established. So no isomorphism between them exists, ever, no matter how the map is defined.

It is precisely $\mathbb{Z} \not\cong \mathbb{Q}$. Injection, yes. Isomorphism, no --- and the elements it misses are the ones you just characterised as $b \neq 1$.

\textbf{Your move.} Write the corrected arrow for $\phi$, and add one line saying which of the two clauses above the old arrow failed.


\needspace{3\baselineskip}\subsection*{Not quite -- You shrank the target. Both ends of 3.2 are fixed.}
You wrote

$$\phi : R[x_1,\ldots,x_n] \longrightarrow \{\, a/b \in F(x_1,\ldots,x_n) \;:\; b = 1_F \,\}.$$

That is a real technique --- cut the codomain down to the image and a non-surjective map becomes surjective. Two things are wrong with it here.

\textbf{It is not the image.} Take $R = \mathbb{Z}$, so $F = \mathbb{Q}$. The element $\tfrac{1}{2}x_1$ can be written with denominator $1_F$, so it is in your set. Nothing in $\mathbb{Z}[x_1,x_2]$ maps to it --- the coefficient is not an integer. Your set is a copy of $F[x_1,\ldots,x_n]$, not of $R[x_1,\ldots,x_n]$.

\textbf{And the repair is pointed the wrong way.} 3.2 names \emph{both} objects and asserts they are identified:

\begin{quote}
the field of fractions of $R[x_1,\ldots,x_n]$ can be identified naturally with $F(x_1,\ldots,x_n)$.
\end{quote}

$F(x_1,\ldots,x_n)$ is the target, all of it. It is not yours to trim --- trimming it proves something else, and something much smaller: that a ring is isomorphic to a copy of itself. Nobody needed that.

So the end that has to move is the \textbf{domain}, and it has to become a field, because you established last card that a non-field is never isomorphic to a field.

\textbf{Your move.} You already wrote, on this same page, the construction that turns an integral domain into a field: $F = \{r_1/r_2 : r_1, r_2 \in R\}$. $R[x_1,\ldots,x_n]$ is also an integral domain. Apply that same construction to it, and write the resulting domain of $\phi$.


\needspace{3\baselineskip}\subsection*{Aside -- Nearly. Classes of $2$-tuples, not $2$-tuples.}
Answering the pencilled question: \textbf{yes, that is the right shape}, with one correction to the word.

An element of $F(x_1,\ldots,x_n)$ is not a $2$-tuple of polynomials. It is an \emph{equivalence class} of them. Concretely: take ordered pairs $(a,b)$ with $a,b \in F[x_1,\ldots,x_n]$ and $b \neq 0$, and declare

$$(a,b) \sim (c,d) \iff ad = bc.$$

An element of $F(x_1,\ldots,x_n)$ is one whole class under $\sim$, and $a/b$ is just how that class is written.

Why the distinction is not pedantry --- you met it already, at $4/2 = 2/1$. Here it is with polynomials, $R = \mathbb{Z}$, $n = 2$:

$$(x_1,\; x_2) \quad\text{and}\quad (x_1^2,\; x_1x_2)$$

are different pairs, and $x_1 \cdot x_1 x_2 = x_2 \cdot x_1^2$, so they are the \textbf{same element}. One element, infinitely many pairs naming it. A map defined on pairs has to agree on all of them or it is not defined on $F(x_1,\ldots,x_n)$ at all.

And the embedding you have been using reads, in this language: the polynomial $p$ goes to the class of $(p,\, 1_F)$. Polynomials in, classes of pairs out --- so the two ends of your arrow are objects of different kinds, which is exactly the thing you were poking at with the question mark.

\textbf{Your move.} For $\phi$ to be a bijection, an element of its domain must be the same kind of object as an element of its target: a class of pairs of polynomials. Say which ring those polynomials must come from --- and then write the domain that results.


\needspace{3\baselineskip}\subsection*{Aside -- Then I will give it to you. Here is the domain.}
Fair enough. You have circled this five times and that is my fault, not yours --- take it stated.

The domain is **the field of fractions of $R[x_1,\ldots,x_n]$**. Its elements are classes $p/q$ with $p, q \in R[x_1,\ldots,x_n]$ and $q \neq 0$. So:

$$\phi : \operatorname{Frac}\bigl(R[x_1,\ldots,x_n]\bigr) \longrightarrow F(x_1,\ldots,x_n).$$

That is all 3.2's sentence was ever saying. Both ends are now fields, and both ends now hold classes of pairs of polynomials --- the difference is only which ring the polynomials come from: $R$ on the left, $F$ on the right.

Drop it to the smallest case, $R = \mathbb{Z}$, $n = 1$:

\begin{center}\begin{tabular}{ll}\hline
\textbf{} & \textbf{elements are classes $p/q$ where} \\ \hline
$\operatorname{Frac}(\mathbb{Z}[x])$ & $p, q \in \mathbb{Z}[x]$, $q \neq 0$ \\
$\mathbb{Q}(x)$ & $p, q \in \mathbb{Q}[x]$, $q \neq 0$ \\
\hline\end{tabular}\end{center}

Integer coefficients on the left, rational coefficients on the right, and 3.2 claims these two fields are the same. $\dfrac{x^2+1}{3x}$ is plainly in both.

The one that is not obvious is the other direction --- something written with rational coefficients, landing in the left column.

\textbf{Your move.} $\dfrac{\tfrac{1}{2}x + 1}{x}$ is an element of $\mathbb{Q}(x)$. Write it as $p/q$ with $p$ and $q$ both in $\mathbb{Z}[x]$.


\needspace{3\baselineskip}\subsection*{Correct -- $p = x+2$, $q = 2x$. Correct --- and that was the proof in miniature.}
Correct. $\dfrac{\tfrac12 x + 1}{x} = \dfrac{x+2}{2x}$, and the check is the relation you wrote earlier: $(\tfrac12 x + 1)\cdot 2x = x \cdot (x+2)$, both sides $x^2 + 2x$. Same class, and now both entries sit in $\mathbb{Z}[x]$.

The step that carried it: you multiplied numerator and denominator by $2$ --- the denominator of the offending coefficient. Nothing else happened.

That move is not a trick for this example. It is the whole of the hard direction of 3.2. Everything up to now was setting up the two ends of $\phi$; you are now inside the proof.

To see why, notice what you have just shown: an element of $\mathbb{Q}(x)$, written with rational coefficients, \emph{is} an element of $\operatorname{Frac}(\mathbb{Z}[x])$ --- it was only the choice of representative that made it look otherwise. That is surjectivity of $\phi$, in the case $R = \mathbb{Z}$, $n = 1$, for one element.

So the next thing is to do it for every element, and the only question is whether the multiplier always exists.

\textbf{Your move.} Take an arbitrary $a/b \in \mathbb{Q}(x)$, with $a, b \in \mathbb{Q}[x]$ and $b \neq 0$. Name the integer $d$ you would multiply numerator and denominator by --- and say why such a $d$ always exists, whatever $a$ and $b$ are.


\needspace{3\baselineskip}\subsection*{Aside -- You are not looking for $ab^{-1}$. You are looking for a different representative.}
You wrote: *all I know is that $b^{-1}$ exists but $ab^{-1}$ need not be in $\mathbb{Z}[x]$.* That is true, and it is also not what is being asked --- which is exactly why you are stuck.

$ab^{-1}$ is not the target. An element of $\mathbb{Q}(x)$ \textbf{is} the class $a/b$, and the class does not change when you multiply both entries by the same non-zero $d$: the pairs $(a,b)$ and $(da, db)$ are related, since $a \cdot db = b \cdot da$. So $a/b = (da)/(db)$, same element, different representative.

That is precisely what you did last card. You did not invert $x$. You replaced the pair $(\tfrac12 x + 1,\; x)$ with the pair $(x+2,\; 2x)$ --- same class, and both entries now in $\mathbb{Z}[x]$.

So $d$ has no obligation to simplify anything. Its one job: make $da$ and $db$ \textbf{both} land in $\mathbb{Z}[x]$.

Before the general $a$ and $b$, one concrete pair.

\textbf{Your move.} Take

$$a = \tfrac{1}{2}x^{2} + \tfrac{2}{3}x, \qquad b = \tfrac{5}{6}x + 3 .$$

Name the single integer $d$ for which $da$ and $db$ have \emph{all} their coefficients in $\mathbb{Z}$, and write out $da$ and $db$.


\needspace{3\baselineskip}\subsection*{Not quite -- $d = 6$ is right. "The maximal $q$" is not the rule.}
$d = 6$ is correct. $6a = 3x^{2} + 4x$ and $6b = 5x + 18$, both in $\mathbb{Z}[x]$. That part is settled.

The rule you attached to it is not. The denominators occurring in $a$ and $b$ are $2$, $3$ and $6$, and the largest of them is $6$ --- which happens to be the same number that works. On this example the right answer and your rule coincide, and that is the example's doing, not the rule's.

So test it on a pair where they come apart.

\textbf{Your move.} Take

$$a = \tfrac{1}{2}x + \tfrac{1}{3}, \qquad b = 1 .$$

Your rule says: take the largest $q$ among the coefficients $p/q$. Name the $d$ that gives, write out $da$, and say whether $da \in \mathbb{Z}[x]$.


\needspace{3\baselineskip}\subsection*{Correct -- The lcm. Correct --- and now the word that makes it a proof.}
Correct. The denominators occurring in $a$ and $b$ are $2, 3, 6, 1$, and their lcm is $6$. And on the pair I handed you, $\tfrac12 x + \tfrac13$, the largest denominator is $3$ while the lcm is $6$ --- which is exactly where the old rule would have failed.

So: $d = \operatorname{lcm}$ of the denominators of all coefficients of $a$ and of $b$.

One thing stands between that and a proof. $\operatorname{lcm}$ of a set of positive integers is a \emph{number} only when the set is \textbf{finite} --- an infinite set of denominators has no common multiple to take. Your rule therefore rests on there being finitely many coefficients to look at.

And the margin note at the top of your page says the opposite: a polynomial is a countably infinite tuple of elements of $R$, *"infinitely many of which may be $0_R$"*.

\textbf{Your move.} Write $a = \tfrac{1}{2}x^{2} + \tfrac{2}{3}x$ out as that tuple --- its coefficients in order, starting with the constant term. How many of its entries are non-zero?


\needspace{3\baselineskip}\subsection*{Your move -- Two. So the note is backwards --- and now the arbitrary $a/b$.}
Two: the tuple is $(0, \tfrac23, \tfrac12, 0, 0, \ldots)$. And that is not special to this polynomial. Every polynomial has finitely many non-zero coefficients and \emph{all but finitely many} entries equal to $0_R$ --- which is the reverse of what your margin note says, and you have just shown it yourself. Fix the note; it is the hypothesis your $d$ depends on.

Now the thing the last three cards were for. Nothing new is needed.

\textbf{Your move.} Let $a/b \in \mathbb{Q}(x)$ be arbitrary, with $a, b \in \mathbb{Q}[x]$ and $b \neq 0$. Produce an integer $d$ such that $da$ and $db$ both lie in $\mathbb{Z}[x]$, say why such a $d$ exists for \emph{every} $a$ and $b$, and write down the resulting pair in $\mathbb{Z}[x]$ representing the same element of $\mathbb{Q}(x)$.

What the statement uses:

\begin{itemize}
  \item $\mathbb{Q}(x) = \operatorname{Frac}(\mathbb{Q}[x])$: its elements are classes of pairs $(a,b)$ with $a, b \in \mathbb{Q}[x]$, $b \neq 0$, written $a/b$.
  \item $a/b = c/e$ exactly when $ae = bc$.
  \item Hence $a/b = (da)/(db)$ for any $d \neq 0$: same class, different representative.
  \item A polynomial in $\mathbb{Q}[x]$ has finitely many non-zero coefficients, each of the form $p/q$ with $p, q \in \mathbb{Z}$, $q \neq 0$.
\end{itemize}


\begin{studentturn}{You wrote -- attempt 1 of 2, 2026-09-02 15:48:45}
\begin{center}\includegraphics[width=0.95\linewidth,height=0.78\textheight,keepaspectratio]{live/answers/t0008-r1.png}\end{center}
\end{studentturn}

\begin{studentturn}{You wrote -- attempt 2 of 2, 2026-09-02 15:55:21}
\begin{center}\includegraphics[width=0.95\linewidth,height=0.78\textheight,keepaspectratio]{live/answers/t0008-r2.png}\end{center}
\end{studentturn}

\needspace{3\baselineskip}\subsection*{Aside -- Both questions are right to ask. The lcm was never the point.}
**1. The lcm is $\mathbb{Z}$'s convenience, and you do not need it.** You are right that a general integral domain $R$ has no lcm to appeal to --- no ordering, no factorisation. But look at what $d$ was ever asked to do: clear finitely many denominators at once. \emph{Any} common multiple does that; the least one was never required. And every commutative ring has one available --- the \textbf{product} of the denominators. In $\mathbb{Z}$, for $\tfrac12 x + \tfrac13$, the product $2 \cdot 3 = 6$ clears them, and so would $12$, or $6000$. Minimality is a habit carried over from arithmetic, not a hypothesis of this problem. So the $\mathbb{Q}(x)$ case was not a special trick: replace "lcm" by "product" and the argument stops mentioning $\mathbb{Z}$ at all.

\textbf{2. Finitely many non-zero coefficients is not a requirement --- it is the definition.} A polynomial \emph{is} a tuple with only finitely many non-zero entries. Drop that and you have a formal power series, a different object living in $F[[x_1,\ldots,x_n]]$; without it $F[x_1,\ldots,x_n]$ is not even closed under the multiplication you know, since each coefficient of a product is a \emph{finite} sum. So it is not a condition to impose on $F[x_1,\ldots,x_n]$ --- it is what makes it that ring rather than the power series ring.

\textbf{Your move.} Let $R$ be an integral domain, $F = \operatorname{Frac}(R)$, and suppose a polynomial in $F[x]$ has exactly two non-zero coefficients, $a_1/b_1$ and $a_2/b_2$, with $a_i, b_i \in R$ and $b_i \neq 0$. Name a single $d \in R$, $d \neq 0$, such that multiplying that polynomial by $d$ puts both of its coefficients in the image of $R$ inside $F$.


\needspace{3\baselineskip}\subsection*{Your move -- $d = b_1 b_2$. Correct --- now do it to an arbitrary element.}
Correct, and note where the hypothesis paid for itself: $d = b_1 b_2 \neq 0$ because $R$ is an \textbf{integral domain}. In a ring with zero divisors the product of two non-zero denominators can be $0$, and the whole move collapses. That is what "integral domain" is doing in this problem.

Nothing else is missing. Two coefficients becomes finitely many the same way, and one polynomial becomes two.

\textbf{Your move.} Prove that $\phi$ is surjective: every element of $F(x_1,\ldots,x_n)$ is $\phi$ of something in $\operatorname{Frac}(R[x_1,\ldots,x_n])$.

What the statement uses:

\begin{itemize}
  \item $R$ is an integral domain and $F = \operatorname{Frac}(R)$.
  \item $\operatorname{Frac}(S)$, for an integral domain $S$: classes of pairs $(p,q)$ with $p, q \in S$, $q \neq 0$, written $p/q$, where $p/q = p'/q'$ exactly when $pq' = p'q$.
  \item Consequently $p/q = (dp)/(dq)$ for any $d \neq 0$.
  \item $R[x_1,\ldots,x_n]$: polynomials in $n$ variables with coefficients in $R$ --- finitely many non-zero coefficients, by definition.
  \item $F(x_1,\ldots,x_n) = \operatorname{Frac}(F[x_1,\ldots,x_n])$.
  \item $R \subseteq F$ and hence $R[x_1,\ldots,x_n] \subseteq F[x_1,\ldots,x_n]$, so a pair over $R[x_1,\ldots,x_n]$ is also a pair over $F[x_1,\ldots,x_n]$; $\phi$ sends the class of such a pair in $\operatorname{Frac}(R[x_1,\ldots,x_n])$ to its class in $F(x_1,\ldots,x_n)$.
  \item \emph{Surjective:} for every $\beta$ in the codomain there is some $\alpha$ in the domain with $\phi(\alpha) = \beta$.
\end{itemize}


\begin{studentturn}{You wrote -- attempt 1 of 4, 2026-09-02 16:09:47}
\begin{center}\includegraphics[width=0.95\linewidth,height=0.78\textheight,keepaspectratio]{live/answers/t0009-r1.png}\end{center}
\end{studentturn}

\begin{studentturn}{You wrote -- attempt 2 of 4, 2026-09-02 16:17:56}
\begin{center}\includegraphics[width=0.95\linewidth,height=0.78\textheight,keepaspectratio]{live/answers/t0009-r2.png}\end{center}
\end{studentturn}

\begin{studentturn}{You wrote -- attempt 3 of 4, 2026-09-02 16:24:32}
\begin{center}\includegraphics[width=0.95\linewidth,height=0.78\textheight,keepaspectratio]{live/answers/t0009-r3.png}\end{center}
\end{studentturn}

\begin{studentturn}{You said -- attempt 4 of 4, 2026-09-02 16:31:13}
I have lost patience. Next question time, actually skip 3.3 as well
\end{studentturn}

\needspace{3\baselineskip}\subsection*{Aside -- Yes --- and it never uses cancellation.}
Your question first. **Yes, $d\alpha$ has all its coefficients in the image of $R$**, and the suspicion is worth having --- but nothing in this step uses cancellation.

Watch one coefficient. Write $d = bc$, where $b$ is the denominator of the coefficient $a/b$ and $c$ is the product of all the other denominators. Then

$$d \cdot \frac{a}{b} \;=\; \frac{d}{1}\cdot\frac{a}{b} \;=\; \frac{bca}{b}.$$

Whether that equals $\dfrac{ca}{1}$ is not a question of cancelling $b$ from top and bottom --- you are right that you have no such licence in a general $R$. It is a question of the \emph{defining relation} of $\operatorname{Frac}(R)$: $p/q = p'/q'$ exactly when $pq' = p'q$. You are not permitted to cancel; you are permitted to cross-multiply, and cross-multiplying is the only thing this argument ever needed. It is the same move you made on $\dfrac{\tfrac12 x + 1}{x} = \dfrac{x+2}{2x}$ hours ago.

One correction to your setup, and it is the other half of the same worry: **$R$ is not an arbitrary ring.** It is an integral domain. Without that, $\operatorname{Frac}(R)$ does not exist to write down, and $d$ --- a product of non-zero denominators --- could be $0$.

\textbf{Your move.} Let $R$ be an integral domain and $a, b, c \in R$ with $b \neq 0$, $c \neq 0$. Show that

$$\frac{bca}{b} \;=\; \frac{ca}{1}$$

holds in $F = \operatorname{Frac}(R)$: write out the cross-multiplication the definition demands, and say why the two sides are equal.


\needspace{3\baselineskip}\subsection*{Aside -- Mostly. But you are not allowed to write $\phi^{-1}$ yet.}
You asked, so here is the audit.

\textbf{Right.} $R$ an integral domain --- fixed. $F$ as you describe it. $d$ = the product of the denominators of the non-zero coefficients of $\alpha$ and $\beta$: right, and it is the only $d$ you need. $d\alpha, d\beta \in R[x]$: true --- though that is the claim I asked you to check by cross-multiplication last card, and you have asserted it instead.

\textbf{Not right.} $\phi^{-1}$ presumes $\phi$ is a bijection, and that is what you are trying to prove. You do not have an inverse. You have a \textbf{candidate preimage}, and surjectivity is exactly this: given $\alpha/\beta$, \emph{exhibit} an element of $\operatorname{Frac}(R[x])$, then \emph{check} that $\phi$ sends it to $\alpha/\beta$. Nothing more, and nothing you can skip.

So the line is not $\phi^{-1}(\alpha/\beta) = d\alpha/d\beta$. It is: set $\gamma = d\alpha/d\beta$, an element of $\operatorname{Frac}(R[x])$ --- and then the check. (It is an element only if $d\beta \neq 0$; that clause belongs in your proof, with its reason.)

The check is the part you have not written at all, and it is one line.

\textbf{Your move.} With $\gamma = d\alpha/d\beta \in \operatorname{Frac}(R[x])$, verify that $\phi(\gamma) = \alpha/\beta$ in $F(x)$.

\begin{itemize}
  \item $\phi$ sends the class of the pair $(p, q)$ over $R[x]$ to the class of the same pair over $F[x]$.
  \item Two classes are equal in $F(x)$ exactly when the cross-products agree: $p/q = p'/q'$ iff $pq' = p'q$.
\end{itemize}


\needspace{3\baselineskip}\subsection*{Aside -- Stuck is fine. Here is what "$d\alpha \in R[x]$" is actually asking.}
I think the trouble is the \emph{target}, not the algebra. Inside $F$, "this element is in $R$" cannot mean literal containment --- elements of $F$ are classes of pairs. It means: **the class equals $r/1$ for some $r \in R$.** That is the shape you are trying to reach, and once you know it the work is one cross-multiplication.

Two rules, both already yours:

\begin{itemize}
  \item Multiplication in $F$: $\dfrac{p}{q}\cdot\dfrac{p'}{q'} = \dfrac{pp'}{qq'}$, and $d \in R$ sits in $F$ as $d/1$.
  \item Equality in $F$: $\dfrac{p}{q} = \dfrac{p'}{q'}$ exactly when $pq' = p'q$.
\end{itemize}

Concretely first, in $R = \mathbb{Z}$. Say the coefficient is $2/3$ and the only other denominator is $5$, so $d = 15$. Then

$$\frac{15}{1}\cdot\frac{2}{3} = \frac{30}{3} = \frac{10}{1},$$

and that last equality is not cancellation --- it is $30 \cdot 1 = 3 \cdot 10$. The integer produced is $10 = 5 \cdot 2$: the \emph{other} denominator times the numerator.

Now the same thing with letters. Fix one non-zero coefficient $a/b$ of $\alpha$. Since $d$ is the product of \textbf{all} the denominators, $b$ is one of its factors: write $d = bc$, where $c \in R$ is the product of the rest. Then

$$\frac{d}{1}\cdot\frac{a}{b} = \frac{bca}{b}.$$

\textbf{Your move.} Name the element $r \in R$ for which $\dfrac{bca}{b} = \dfrac{r}{1}$, and check that equality by writing out the cross-multiplication the definition demands.

That single line, said for an arbitrary coefficient, is the whole of $d\alpha \in R[x]$ --- and $d\beta$ goes the same way, since $d$ was built from $\beta$'s denominators too.


\needspace{3\baselineskip}\subsection*{Your move -- 3.8. First, two values in $\mathbb{F}_2$.}
Next on the sheet is \textbf{3.8}, in the book's words:

\begin{quote}
Suppose that $K$ is a field. If $f = a_0 + a_1x + \cdots + a_nx^n \in K[x]$ and $k \in K$, let $(\Phi(f))(k) = a_0 + a_1k + \cdots + a_nk^n$. Show that $\Phi$ is a ring homomorphism from $K[x]$ to $K^K$. Show that if $K$ is finite then $\Phi$ is an epimorphism, but not a monomorphism. What happens if $K$ is infinite?
\end{quote}

You are not answering that yet.

What the statement uses:

\begin{itemize}
  \item $K^K$: the set of \textbf{all} functions from $K$ to $K$, made a ring by working pointwise --- $(g+h)(k) = g(k)+h(k)$ and $(gh)(k) = g(k)h(k)$, with zero the constant function $0$ and identity the constant function $1$.
  \item $\Phi(f)$ is therefore not a number: it is the \emph{function} sending each $k \in K$ to $f$ evaluated at $k$.
  \item \emph{Epimorphism:} a surjective homomorphism. \emph{Monomorphism:} an injective one.
  \item $\mathbb{F}_2 = \{0, 1\}$, the field with two elements, where $1 + 1 = 0$.
\end{itemize}

\textbf{Your move.} Take $K = \mathbb{F}_2$ and $f = x^2 + x \in \mathbb{F}_2[x]$. Write down the function $\Phi(f)$ by giving both of its values: $(\Phi(f))(0)$ and $(\Phi(f))(1)$.


\begin{studentturn}{You wrote -- attempt 1 of 15, 2026-09-02 16:39:46}
\begin{center}\includegraphics[width=0.95\linewidth,height=0.78\textheight,keepaspectratio]{live/answers/t0011-r1.png}\end{center}
\end{studentturn}

\begin{studentturn}{You wrote -- attempt 2 of 15, 2026-09-02 16:44:26}
\begin{center}\includegraphics[width=0.95\linewidth,height=0.78\textheight,keepaspectratio]{live/answers/t0011-r2.png}\end{center}
\end{studentturn}

\begin{studentturn}{You wrote -- attempt 3 of 15, 2026-09-02 17:59:44}
\begin{center}\includegraphics[width=0.95\linewidth,height=0.78\textheight,keepaspectratio]{live/answers/t0011-r3.png}\end{center}
\end{studentturn}

\begin{studentturn}{You wrote -- attempt 4 of 15, 2026-09-02 18:19:13}
\begin{center}\includegraphics[width=0.95\linewidth,height=0.78\textheight,keepaspectratio]{live/answers/t0011-r4.png}\end{center}
\end{studentturn}

\begin{studentturn}{You wrote -- attempt 5 of 15, 2026-09-02 18:26:03}
\begin{center}\includegraphics[width=0.95\linewidth,height=0.78\textheight,keepaspectratio]{live/answers/t0011-r5.png}\end{center}
\end{studentturn}

\begin{studentturn}{You wrote -- attempt 6 of 15, 2026-09-02 18:35:48}
\begin{center}\includegraphics[width=0.95\linewidth,height=0.78\textheight,keepaspectratio]{live/answers/t0011-r6.png}\end{center}
\end{studentturn}

\begin{studentturn}{You wrote -- attempt 7 of 15, 2026-09-02 18:42:26}
\begin{center}\includegraphics[width=0.95\linewidth,height=0.78\textheight,keepaspectratio]{live/answers/t0011-r7.png}\end{center}
\end{studentturn}

\begin{studentturn}{You wrote -- attempt 8 of 15, 2026-09-02 18:46:45}
\begin{center}\includegraphics[width=0.95\linewidth,height=0.78\textheight,keepaspectratio]{live/answers/t0011-r8.png}\end{center}
\end{studentturn}

\begin{studentturn}{You wrote -- attempt 9 of 15, 2026-09-02 18:46:48}
\begin{center}\includegraphics[width=0.95\linewidth,height=0.78\textheight,keepaspectratio]{live/answers/t0011-r9.png}\end{center}
\end{studentturn}

\begin{studentturn}{You wrote -- attempt 10 of 15, 2026-09-02 18:57:25}
\begin{center}\includegraphics[width=0.95\linewidth,height=0.78\textheight,keepaspectratio]{live/answers/t0011-r10.png}\end{center}
\end{studentturn}

\begin{studentturn}{You wrote -- attempt 11 of 15, 2026-09-02 18:57:33}
\begin{center}\includegraphics[width=0.95\linewidth,height=0.78\textheight,keepaspectratio]{live/answers/t0011-r11.png}\end{center}
\end{studentturn}

\begin{studentturn}{You wrote -- attempt 12 of 15, 2026-09-02 19:07:34}
\begin{center}\includegraphics[width=0.95\linewidth,height=0.78\textheight,keepaspectratio]{live/answers/t0011-r12.png}\end{center}
\end{studentturn}

\begin{studentturn}{You wrote -- attempt 13 of 15, 2026-09-02 19:29:55}
\begin{center}\includegraphics[width=0.95\linewidth,height=0.78\textheight,keepaspectratio]{live/answers/t0011-r13.png}\end{center}
\end{studentturn}

\begin{studentturn}{You wrote -- attempt 14 of 15, 2026-09-02 19:36:04}
\begin{center}\includegraphics[width=0.95\linewidth,height=0.78\textheight,keepaspectratio]{live/answers/t0011-r14.png}\end{center}
\end{studentturn}

\begin{studentturn}{You wrote -- attempt 15 of 15, 2026-09-02 19:44:10}
\begin{center}\includegraphics[width=0.95\linewidth,height=0.78\textheight,keepaspectratio]{live/answers/t0011-r15.png}\end{center}
\end{studentturn}

\needspace{3\baselineskip}\subsection*{Aside -- The value lives in $K$. The function lives in $K^K$.}
The page shows exactly where you are lost, and it is one thing.

You wrote $(\Phi(f))(0) = 0 \in K^K$. That $0$ is \textbf{not} in $K^K$. It is in $K$.

\begin{itemize}
  \item $\Phi(f)$ --- on its own, with nothing fed to it --- is the element of $K^K$. It is a \emph{function}.
  \item $(\Phi(f))(0)$ --- that function handed the input $0$ --- is a single element of $K$.
\end{itemize}

$K = \mathbb{F}_2$ has two elements, so the function $\Phi(f)$ is completely described by two entries:

\begin{center}\begin{tabular}{ll}\hline
\textbf{$k$} & \textbf{$(\Phi(f))(k)$} \\ \hline
$0$ & $0$ \\
$1$ & ? \\
\hline\end{tabular}\end{center}

Your first line filled the top row and filled it correctly: $0^2 + 0 = 0$. Nothing else about that line was wrong except where you said the answer lived.

Arithmetic you may want, and it is not what this exercise is about: in $\mathbb{F}_2$, $1 \cdot 1 = 1$ and $1 + 1 = 0$.

\textbf{Your move.} Fill the bottom row: compute $(\Phi(f))(1) = 1^2 + 1$, as an element of $K$.

Once that row is filled you will be able to read the whole function off the table, and one half of 3.8 falls straight out of it.


\needspace{3\baselineskip}\subsection*{Aside -- Both values right. But $\mathbb{F}_2[x] \not\subset K^K$.}
Both values are correct, and the working you showed is the right working: $1 \cdot 1 + 1 = 1 + 1 = 0$. So the table is

\begin{center}\begin{tabular}{ll}\hline
\textbf{$k$} & \textbf{$(\Phi(f))(k)$} \\ \hline
$0$ & $0$ \\
$1$ & $0$ \\
\hline\end{tabular}\end{center}

and $\Phi(f)$ is the \textbf{zero function} --- the additive identity of $K^K$.

One line on the page is false, and it is the first: you wrote $f = x^2 + x \in \mathbb{F}_2[x] \subset K^K$. That containment does not hold. A polynomial is not a function --- you established that yourself earlier today --- and $\Phi$ is a \emph{map from} $K[x]$ \emph{to} $K^K$, not an inclusion of one inside the other. The whole of 3.8 is a measurement of how badly $\Phi$ fails to be an inclusion.

And your own table is the evidence.

\textbf{Your move.} Name a polynomial $g \in \mathbb{F}_2[x]$, different from $f$, with $\Phi(g) = \Phi(f)$ --- the same function, both values agreeing. The table you just filled hands you one immediately.


\needspace{3\baselineskip}\subsection*{Correct -- Correct --- and your $g$ is the zero polynomial}
Correct. $\Phi(g) = \Phi(f)$ as elements of $K^K$ (both send $0 \mapsto 0$ and
$1 \mapsto 0$), and $g \neq f$ in $\mathbb{F}_2[x]$. That is a non-zero kernel,
so **$\Phi$ is not a monomorphism** --- you have just answered that clause of 3.8
for $K = \mathbb{F}_2$.

One thing to notice about what you picked: $x + x = (1+1)x = 0 \cdot x = 0$ in
$\mathbb{F}_2[x]$, so your $g$ \emph{is} the zero polynomial. It still does the job ---
$0 \neq x^2 + x$, and $\Phi(0)$ is the zero function --- but say it as $g = 0$ and
the point is sharper: $x^2 + x$ lies in $\ker \Phi$.

Now the first clause of the problem, which is the one everything else rests on.

\textbf{3.8.} Suppose that $K$ is a field. If $f = a\_0 + a\_1x + \cdots + a\_nx^n \in
K[x]$ and $k \in K$, let $(\Phi(f))(k) = a\_0 + a\_1k + \cdots + a\_nk^n$. Show
that $\Phi$ is a ring homomorphism from $K[x]$ to $K^K$.

What that leans on:

\begin{itemize}
  \item *$K^K$:* all mappings $K \to K$, with $(u+v)(k) = u(k)+v(k)$ and
\end{itemize}
  $(uv)(k) = u(k)v(k)$ --- the ring you built in 3.1.
\begin{itemize}
  \item *Sum in $K[x]$:* coefficientwise, $(f+g)$ has $i$-th coefficient $a_i + b_i$.
  \item \emph{Ring homomorphism:} $\Phi(f+g) = \Phi(f) + \Phi(g)$ and
\end{itemize}
  $\Phi(fg) = \Phi(f)\Phi(g)$.
\begin{itemize}
  \item \emph{Two functions are equal} when they agree at every point of $K$.
\end{itemize}

\textbf{Your move.} The additive half only: let $f = \sum a_ix^i$ and
$g = \sum b_ix^i$ in $K[x]$, and prove $\Phi(f+g) = \Phi(f) + \Phi(g)$. Fix an
arbitrary $k \in K$ and compute both sides at $k$. The product half is the next
turn.


\needspace{3\baselineskip}\subsection*{Aside -- Not by definition --- those are two \emph{different} definitions}
Right question, and you are not underthinking it. But look at what you wrote:
the line

$$(\Phi(f+g))(k) = (\Phi(f))(k) + (\Phi(g))(k)$$

\emph{is} the thing to be proved, at the point $k$. Writing it down asserts it.

It is not immediate, because the two sides are reached by two different
definitions and neither one mentions the other.

\textbf{Left side.} $f+g$ is a single polynomial in $K[x]$, whose $i$-th coefficient
is $a_i + b_i$. Apply the definition of $\Phi$ to \emph{that} polynomial:
$$(\Phi(f+g))(k) = \sum_i (a_i + b_i)k^i.$$

\textbf{Right side.} $\Phi(f) + \Phi(g)$ is a sum in $K^K$, and addition there is
pointwise --- the definition from 3.1. So
$$(\Phi(f) + \Phi(g))(k) = (\Phi(f))(k) + (\Phi(g))(k) = \sum_i a_ik^i + \sum_i b_ik^i.$$

So everything the proof has to do is:
$$\sum_i (a_i + b_i)k^i = \sum_i a_ik^i + \sum_i b_ik^i,$$
an identity in $K$. Small --- one axiom of $K$ applied to each term, then
regrouping a finite sum. Small is not the same as free. That is the only place
in the additive half where anything is actually shown.

(If $f$ and $g$ have different degrees, run both sums to the same $n$ and pad
the shorter one with zero coefficients. That is what $f+g$ means anyway.)

\textbf{Your move.} Rewrite it as three lines: left side expanded by the definition
of $\Phi$; right side expanded by pointwise addition and then the definition of
$\Phi$; and the equality joining them. Name the axiom of $K$ that turns
$(a_i + b_i)k^i$ into $a_ik^i + b_ik^i$.


\needspace{3\baselineskip}\subsection*{Correct -- Correct --- the additive half is done. Now the product.}
That is the proof. Every equality is an application of exactly one definition,
in the right order: $\Phi$ applied to the coefficientwise sum, then $\Phi$'s own
formula, then the identity in $K$, then back out again. And $k$ was arbitrary,
so the two functions agree everywhere and are equal in $K^K$.

The step you did not label --- line 3 to line 4 --- is \textbf{distributivity} in $K$,
$(a_i + b_i)k^i = a_ik^i + b_ik^i$ termwise, then commutativity and
associativity of $+$ to regroup the finite sum into two. Write the word next
time; it is the only axiom the additive half uses.

Now the other half, and this one has real arithmetic in it.

**Show that $\Phi(fg) = \Phi(f)\Phi(g)$.**

What it leans on:

\begin{itemize}
  \item *Product in $K[x]$:* $fg$ is the polynomial whose $m$-th coefficient is
\end{itemize}
  $\sum_{i+j=m} a_ib_j$ --- that is, $c_m = a_0b_m + a_1b_{m-1} + \cdots + a_mb_0$.
\begin{itemize}
  \item *Product in $K^K$:* pointwise, $(uv)(k) = u(k)\,v(k)$.
  \item *$\Phi$:* $(\Phi(f))(k) = \sum_i a_ik^i$.
  \item \emph{Equal functions:} they agree at every $k \in K$.
\end{itemize}

\textbf{Your move.} Same shape as what you just wrote. Fix an arbitrary $k \in K$,
expand $(\Phi(fg))(k)$ using the two definitions above, expand
$(\Phi(f)\Phi(g))(k)$ using its two, and get from one to the other.

If you want the arithmetic in your hand before the general case, do it first for
$f = a_0 + a_1x$ and $g = b_0 + b_1x$: multiply them out in $K[x]$, evaluate at
$k$, and multiply $(a_0 + a_1k)(b_0 + b_1k)$ out separately. Compare. Your
choice --- the general statement is what is being asked for.


\needspace{3\baselineskip}\subsection*{Aside -- FOILing is exactly what it is --- collected by power}
Your question first.

**Multiplication in $K[x]$ is FOILing.** That is not a shortcut or a trick; it
\emph{is} the definition, written out. You multiply every term of $f$ by every term
of $g$, and then collect the ones that carry the same power of $x$.

Do it for $f = a_0 + a_1x$ and $g = b_0 + b_1x$ and watch it happen:

$$(a_0 + a_1x)(b_0 + b_1x) = \underbrace{a_0b_0}_{x^0} + \underbrace{(a_0b_1 + a_1b_0)}_{x^1}x + \underbrace{a_1b_1}_{x^2}x^2$$

The middle coefficient is a \emph{sum of two products}, because there were two ways
to make an $x^1$: take the $x^0$ from $f$ and the $x^1$ from $g$, or the other
way round. That is the whole content of the general formula:

$$fg = \sum_m c_m x^m, \qquad c_m = \sum_{i+j=m} a_ib_j.$$

The $m$-th coefficient of $fg$ collects every pair of indices that adds to $m$.
So $c_2 = a_0b_2 + a_1b_1 + a_2b_0$ --- three terms, not one.

What it is \textbf{not}: the $m$-th coefficient of $fg$ is not $a_mb_m$. Coefficients
do not multiply in place the way they add in place. Addition is index-by-index;
multiplication is not, and that asymmetry is the only genuinely new thing in
this half of 3.8.

\textbf{Your move} --- thirty seconds, so the arithmetic is in your hand before the
general case. In $\mathbb{Z}[x]$ take

$$f = 2 + 3x, \qquad g = 1 + 4x + 5x^2.$$

Write down $c_0$, $c_1$, $c_2$, $c_3$ --- the four coefficients of $fg$ --- using
$c_m = \sum_{i+j=m} a_ib_j$. Four numbers, that is all.


\needspace{3\baselineskip}\subsection*{Not quite -- $c_0$ and $c_1$ are right. You stopped collecting pairs after $c_1$.}
$c_0 = 2$. Correct.

$c_1 = 3 + 8 = 11$. Correct, and it is the important one: you took \textbf{both}
ways of making an $x^1$, $a_0b_1$ and $a_1b_0$, and added them.

Then you stopped doing that.

**$c_2 = 10$.** List the pairs $(i,j)$ with $i + j = 2$, where $i$ ranges over
$f$'s indices $\{0,1\}$ and $j$ over $g$'s indices $\{0,1,2\}$. There are two of
them. You have used one.

**$c_3 = 0$.** Same question for $i + j = 3$. There is a pair. Find it, and
then $c_3$ is not $0$.

\medskip\hrule\medskip

And the line above your bubble, which is the same error one level up:

$$(\Phi(fg))(k) = \Big(\Phi\big(\textstyle\sum (a_i b_i)x^i\big)\Big)(k)$$

$\sum (a_ib_i)x^i$ is \textbf{not} $fg$. Your own $c_1$ says so: you computed
$a_0b_1 + a_1b_0$, not $a_1b_1$. The coefficient of $x^m$ in $fg$ is
$c_m = \sum_{i+j=m} a_ib_j$ --- the whole sum, not the single term with $i = j = m$.
Rewrite that first line with $c_m$ in it.

One notation point while you are there: in $K^K$ the product $\Phi(f)\Phi(g)$ is
the \textbf{pointwise} product, $(\Phi(f)\Phi(g))(k) = \Phi(f)(k)\cdot\Phi(g)(k)$ ---
multiply the two values in $K$. It is not composition, so do not write $\circ$
for it; that would be a different claim.

\textbf{Your move.} Fix $c_2$ and $c_3$, then rewrite that first line.


\needspace{3\baselineskip}\subsection*{Aside -- That is the same page twice --- the card above is its answer}
Both sends carry identical ink, three seconds apart. Nothing new to read, so I
am not marking it a second time: \textbf{card 0039 above is the reply to it.}

The ask still standing, so it is at the bottom of the board where you are
looking rather than a card up:

\begin{itemize}
  \item **$c_2$** --- the pairs $(i,j)$ with $i+j=2$, $i \in \{0,1\}$, $j \in \{0,1,2\}$.
\end{itemize}
  There are two. You used one.
\begin{itemize}
  \item **$c_3$** --- same question for $i+j=3$. There is one, and it is not nothing.
  \item \textbf{The first line of your proof} --- rewrite it with $c_m = \sum_{i+j=m} a_ib_j$
\end{itemize}
  in place of $\sum(a_ib_i)x^i$.


\needspace{3\baselineskip}\subsection*{Aside -- Arithmetic settled. The formula is in the margin, not in the proof.}
$c_2 = 10 + 12 = 22$ and $c_3 = 15$. Both right, and $\sum_{i+j=m} a_ib_j$ is
written out correctly. The rung is done --- the mechanics of multiplying in $K[x]$
are yours.

Now look at where you put it. The arrow runs from $c_m = \sum_{i+j=m} a_ib_j$
into the \textbf{bubble} --- into the $2+3x$ example. The line above the bubble, the one
that is actually your proof, is untouched:

$$(\Phi(fg))(k) = \Big(\Phi\big(\textstyle\sum (a_i b_i)x^i\big)\Big)(k)$$

Those two statements contradict each other. Your formula says the coefficient of
$x^m$ in $fg$ is a sum over every pair with $i+j=m$. Your proof line says it is
the single product $a_mb_m$. In the example you just did, the first gives
$c_2 = 22$ and the second gives $a_2b_2$, which is $0 \cdot 5 = 0$. They cannot
both be $fg$, and the one in the margin is the correct one.

So the substitution that belongs on that line is $fg = \sum_m c_m x^m$ --- the
thing you are applying $\Phi$ to is $fg$, and \emph{that} is what $fg$ is.

\textbf{Your move.} Cross out that line and write it again, with $c_m$ in it. Then
carry on: apply $\Phi$ and evaluate at $k$.


\needspace{3\baselineskip}\subsection*{Aside -- Identical ink again --- 0041 stands}
Same page as the last one, byte for byte, eight seconds later. That is the
second double-send tonight; one tap on \textbf{Send} is enough, and a second one
costs you a turn waiting for me to tell you nothing.

Card 0041 above is the reply. Still owed, unchanged:

**Cross out $\big(\Phi(\sum (a_ib_i)x^i)\big)(k)$ and write that line again with
$c_m$ in it** --- $fg = \sum_m c_m x^m$ is what you are applying $\Phi$ to. Then
apply $\Phi$ and evaluate at $k$.


\needspace{3\baselineskip}\subsection*{Correct -- Correct. The product half is done --- now the epimorphism.}
That is right. The line that carried it is the third one: from
$\sum_m \bigl(\sum_{i+j=m} a_i b_j\bigr) k^m$ back out to
$\bigl(\sum_i a_i k^i\bigr)\bigl(\sum_j b_j k^j\bigr)$ --- the same collecting by
total degree you did with $x$, read backwards, with $k$ in its place. Both
halves of the homomorphism clause are now settled, and I am typing it up.

\textbf{Next piece of 3.8, in the book's words:}

\begin{quote}
Show that if $K$ is finite then $\Phi$ is an \emph{epimorphism}.
\end{quote}

Not yet --- one thing first.

\begin{itemize}
  \item *$K^K$:* every mapping from $K$ to $K$, with the operations of 3.1.
  \item \emph{Epimorphism:} a ring homomorphism that is onto. So: every $h \in K^K$ is
\end{itemize}
  $\Phi(f)$ for some $f \in K[x]$.

\medskip\hrule\medskip

\textbf{Your move.} Take $K = \Fq{2} = \{0, 1\}$. Let $h \in \Fq{2}^{\Fq{2}}$ be the
mapping
$$h(0) = 1, \qquad h(1) = 0.$$

Write down a polynomial $f \in \Fq{2}[x]$ with $\Phi(f) = h$ --- that is, one with
$\bigl(\Phi(f)\bigr)(0) = 1$ and $\bigl(\Phi(f)\bigr)(1) = 0$. Show both
evaluations.

Remember $1 + 1 = 0$ in $\Fq{2}$.


\needspace{3\baselineskip}\subsection*{Correct -- $f = 1+x$. Correct, both evaluations shown.}
Right. $(\Phi(f))(0) = 1+0 = 1$ and $(\Phi(f))(1) = 1+1 = 0$, so $\Phi(1+x) = h$.
($x^2+1$ would have done it too.)

What you just built is the whole finite-$K$ proof in miniature: a polynomial
whose \emph{values} were prescribed at every point of $K$. Over $\Fq{2}$ you could
find it by staring. Over a bigger $K$ you need a recipe, and the recipe is built
out of polynomials that hit one point and kill all the others.

One rung toward that, so the mechanics are fluent before the proof.

\begin{itemize}
  \item *$K$ finite:* write its elements as $c_1,\ldots,c_q$, all distinct.
  \item Recall that in $K[x]$ a product is zero at $k$ exactly when one of its factors
\end{itemize}
  is zero at $k$ --- evaluation is a homomorphism, which is the clause you just
  finished proving.

\medskip\hrule\medskip

\textbf{Your move.} Take $K = \Fq{3} = \{0,1,2\}$, so $1+1 = 2$, $1+1+1 = 0$, and
$2 \cdot 2 = 1$.

Let
$$p(x) = (x-1)(x-2) \in \Fq{3}[x].$$

Compute all three values $p(0)$, $p(1)$, $p(2)$ in $\Fq{3}$. Show the arithmetic
--- note that $-1 = 2$ and $-2 = 1$ here.


\needspace{3\baselineskip}\subsection*{Aside -- Subtraction is not a third operation. It is $+$ with the additive inverse.}
Right to ask, and you are right that the operations are only $+$ and the
product. Subtraction is notation, not a new operation:

\begin{itemize}
  \item A ring's axioms make $(R, +)$ an \textbf{abelian group}. So every $a \in R$ has an
\end{itemize}
  additive inverse $-a$, the unique element with $a + (-a) = 0$.
\begin{itemize}
  \item $a - b$ is \emph{defined} to mean $a + (-b)$. Nothing is being assumed beyond the
\end{itemize}
  additive group you already have.

So in $\Fq{3}$, where $1 + 2 = 0$, the inverse of $1$ is $2$ and the inverse of
$2$ is $1$. Your $p$ has no minus signs in it at all:
$$p(x) = (x-1)(x-2) = (x+2)(x+1) \in \Fq{3}[x].$$

And your three values are correct: $p(0) = 2\cdot 1 = 2$, $p(1) = 0 \cdot 2 = 0$,
$p(2) = 1 \cdot 0 = 0$. One point where it is nonzero, zero at the other two ---
which is the whole reason that product was written down.

\medskip\hrule\medskip

\textbf{Your move.} It is not yet an indicator, because $p(0) = 2$ and you want $1$.
Fix that by scaling.

Find the $c \in \Fq{3}$ with $c \cdot 2 = 1$, set
$$\delta(x) = c \cdot p(x) \in \Fq{3}[x],$$
and give the three values $\delta(0)$, $\delta(1)$, $\delta(2)$.

Recall that in $\Fq{3}$: $2 \cdot 2 = 4 = 1$.


\needspace{3\baselineskip}\subsection*{Correct -- $c = 2$, and $\delta$ is $1,0,0$. Now put three of them together.}
Correct. $\delta(x) = 2(x-1)(x-2)$ takes the values $\delta(0) = 1$,
$\delta(1) = 0$, $\delta(2) = 0$ --- an indicator at $0$.

The same recipe at the other two points, which is your operation run twice more,
so I am not making you do it:

\begin{center}\begin{tabular}{lllll}\hline
\textbf{} & \textbf{polynomial} & \textbf{at $0$} & \textbf{at $1$} & \textbf{at $2$} \\ \hline
$\delta_{0}$ & $2(x-1)(x-2)$ & $1$ & $0$ & $0$ \\
$\delta_{1}$ & $2x(x-2)$ & $0$ & $1$ & $0$ \\
$\delta_{2}$ & $2x(x-1)$ & $0$ & $0$ & $1$ \\
\hline\end{tabular}\end{center}

Each one is built the same way: multiply the two factors that kill the \emph{other}
two points, then scale so the surviving value is $1$.

\medskip\hrule\medskip

\textbf{Your move.} Let $h \in \Fq{3}^{\Fq{3}}$ be the mapping
$$h(0) = 2, \qquad h(1) = 0, \qquad h(2) = 1.$$

Using only $\delta_{0}$, $\delta_{1}$, $\delta_{2}$ from the table and elements
of $\Fq{3}$, write down a polynomial $f \in \Fq{3}[x]$ with $\Phi(f) = h$ --- and
then compute $f(0)$, $f(1)$, $f(2)$ from the table's rows to check that it is.


\end{document}
